% ===================================================================
%  அட்டவணை எண் 11 — நகரமும் அதன் நினைவுச் சின்னங்களும். தீ விபத்து.
%
%  Traduction, faite en 2026, en tamoul d'aujourd'hui, sur l'ido de
%  texto/io. Regles de la colonne : voir 10-attavanai-01.tex. Deux
%  scenes, deux numerotations : les cles portent c1 et c2. Ecrit avec
%  outils/ordo.py sous les yeux, bloc par bloc.
%
%  LES INSTITUTIONS SE DECRIVENT, ELLES NE SE TRANSPOSENT PAS —
%  CINQUIEME LANGUE, CINQUIEME FOIS. Le hindi, le marathe, le
%  telougou et le coreen avaient tous les quatre refuse de donner a
%  la prefecture, a la Bourse et a l'ecole normale de 1926 le nom de
%  l'institution la plus proche de chez eux. Le tamoul avait ici la
%  tentation la plus forte de toutes : le Tamil Nadu a une
%  மாவட்ட ஆட்சியர் அலுவலகம் bien reelle, et le mot serait tombe tout
%  seul sous la plume pour le renvoi (15). Ce n'est pas une
%  prefecture francaise, et l'ecrire aurait deplace la chose de six
%  mille kilometres. La colonne ecrit ஆட்சி மாளிகை, வணிக மாளிகை,
%  ஆசிரியர் பயிற்சிப் பள்ளி — ce que les quatre autres avaient ecrit,
%  chacune pour elle-meme. A cinq rédactions séparées, la convergence
%  n'est plus une coincidence de famille : c'est la regle qui tient.
%
%  ET LE TAMOUL A SES DEUX INSTRUMENTS SUR LES CINQ. L'orchestre du
%  bloc c2-06-1 fuit avec violon, flute, violoncelle, trombone et
%  tambour. Trois s'empruntent — வயலின், செலோ, ட்ரோம்போன் —, deux
%  non : புல்லாங்குழல் est la flute tamoule, et முரசு est le tambour
%  du Sangam, mot vieux de deux mille ans. La coupure ne suit pas
%  l'instrument mais son age dans le pays.
%
%  பால்கனி (77) EST ECRIT ICI ENCORE, comme aux tableaux 5 et 10.
%  Septieme fois que la colonne prefere l'emprunt courant a un mot
%  tamoul vivant mais faux — திண்ணை, மொட்டை மாடம். La regle ne
%  change pas d'un tableau a l'autre, et c'est justement ce qui en
%  fait une regle.
%
%  TROIS MOTS DE POLICE POUR TROIS CORPS DIFFERENTS. Le fac-simile
%  distingue les jendarmi a cheval (61), les jendarmi a pied (63) et
%  les policisti (64) ; le tamoul dirait காவலர் pour les trois et
%  confondrait les renvois. La colonne ecrit குதிரைப் படைக் காவலர்கள்,
%  படைக் காவலர்கள், காவல் துறையினர். Qualifier n'est pas forger : on
%  ne cree pas de mot, on precise avec ceux qui existent — meme parti
%  qu'au தச்சர் du tableau 5.
%
%  LE FAC-SIMILE REEMPLOIE (81) DANS LE MEME BLOC — les musiciens
%  d'abord, le violon ensuite — et la source ne bouge pas.
%
%  VINGT-ET-UNE PAIRES, DOUZE RETOURNEMENTS. La premiere scene
%  enumere les monuments d'une ville, la seconde attache : l'incendie
%  NE dans le foyer, les pompiers QUI viennent secourir, l'affiche
%  QUI annonce, le tuyau DE la pompe. C'est encore le meme partage.
% ===================================================================

%%K t11-apar-1 apar
\VUsaut{2.0mm}
\VUcentre{13.2pt}{90}{மூன்றாம் தொடர்}
\VUsaut{3.0mm}
\VUfilet{16mm}
\VUsaut{5.6mm}
\VUtitre{15.0pt}{60}{அட்டவணை எண் 11}
\VUsaut{1.4mm}
\VUpk{11.4pt}{40}{நகரமும் அதன் நினைவுச் சின்னங்களும். --- தீ விபத்து.}
\VUsaut{2.2mm}

%%K t11-tit-1 sub
\VUpk{10.2pt}{40}{புறநகர்.}
\VUsaut{3.0mm}

%%K t11-c1-01-1 p
1. --- பெருங்கடலிலிருந்து நூறு கிலோமீட்டர் தொலைவில் அமைந்த ஒரு பெரிய
நகரத்தில் நான் வாழ்கிறேன் ; இருந்தாலும் அது உயர்தரமான \VUgras{துறைமுகம்}
\textsuperscript{(1)}. அதன் ஓரமாக ஓடும் ஆற்றின் அகலம் நானூறு மீட்டர் ;
அது மிக அழகான துறையை உண்டாக்குகிறது.

%%K t11-c1-02-1 p
2. --- \VUgras{கற்கரைகளின்} \textsuperscript{(2)} மேல் இருக்கும் நகரப்
பகுதி முழுவதும் கடல் சார்ந்ததாகவும் தொழில் சார்ந்ததாகவும் தெரிகிறது ;
அது ஒரு \VUgras{புறநகரை} \textsuperscript{(3)} உண்டாக்குகிறது ; அங்கே
கப்பலாட்களும் தொழிலாளர்களும் மட்டுமே வாழ்கிறார்கள். இவர்கள்
\VUgras{தொழிற்சாலைகளிலும்} \textsuperscript{(4)} \VUgras{வாயுத்
தொழிற்சாலையிலும்} \textsuperscript{(5)} வேலை செய்கிறார்கள் ; அந்த
வாயுத் தொழிற்சாலையின் உயரமான \VUgras{புகைபோக்கியும்}
\textsuperscript{(6)} \VUgras{வாயுத் தொட்டியும்} வெகு தொலைவிலிருந்தே
தெரிகின்றன.

%%K t11-c1-03-1 p
3. --- இடைக்காலத்தில் நகரம் அரண் செய்யப்பட்டிருந்தது ; அதன் மதில்களின்
சில எச்சங்களை இன்னும் காணலாம் — குறிப்பாக ஒரு \VUgras{வாயில்} \textsuperscript{(8)} ; அது பழைய
\VUgras{அரண் கோட்டையினுடையது} \textsuperscript{(7)}, அதன் இரு பக்கமும் இரண்டு \VUgras{காவல் கோபுரங்கள்}
\textsuperscript{(9)} ; அவை இப்போது \VUgras{சிறைச்சாலையாகப்}
பயன்படுகின்றன.

%%K t11-c1-04-1 p
4. --- மிகவும் பழையதாகத் தெரியும் இந்தப் \VUgras{பகுதி} முழுவதிலும்
ஒரு கட்டிடம் மட்டுமே ஒப்பீட்டளவில் நவீனமானது : அது \VUgras{சுங்க
அலுவலகம்} \textsuperscript{(10)}. அது கற்கரைக்கும் \VUgras{சிறு தெருவுக்கும்}
\textsuperscript{(11)} இடையே அமைந்திருக்கிறது ; அந்தத் தெரு பழைய
\VUgras{நகர மன்ற மாளிகைக்கு} \textsuperscript{(12)} அழைத்துச் செல்கிறது.
அந்தக் கடைசி நினைவுச் சின்னத்தை எளிதாக அடையாளம் காணலாம் ; ஏனெனில்
அதன் மாபெரும் \VUgras{மணிக் கோபுரம்} \textsuperscript{(13)} பழைய
நகரத்தை ஆள்கிறது.

%%K t11-c1-05-1 p
5. --- தெருவின் மறு பக்கத்தில் சுங்க அலுவலகத்தின் அதே பாணியில் கட்டிய
கட்டிடம் ஒன்று \VUgras{நிற்கிறது} : அது \VUgras{வணிக மாளிகை}
\textsuperscript{(14)}. அதன் ஒரு முகப்பு ஒரு சிறிய சதுக்கத்தைப்
பார்க்கிறது ; அங்கே \VUgras{ஆட்சி மாளிகை} \textsuperscript{(15)}
இருக்கிறது. உண்மையில் நம் நகரம் தலைநகரம் அல்லாவிட்டாலும் முக்கியமான
மாவட்டத் தலைநகரம்.

%%K t11-tit-2 sub
\VUsaut{5.0mm}
\VUpk{10.2pt}{40}{நகரத்தின் மிக உயர்ந்த பகுதி.}
\VUsaut{3.4mm}

%%K t11-c1-06-1 p
6. --- படைப் பிரிவு போதுமான அளவு பெரியது ; அது மிகவும் நலம் தரும்
அகன்ற \VUgras{படை முகாம்களில்} \textsuperscript{(16)}
தங்கியிருக்கிறது ; அவை \VUgras{நகரப் பூங்காவின்} \textsuperscript{(17)}
கிராதி வரை நீள்கின்றன. இந்தப் பூங்காவுக்கு அழைத்துச் செல்லும்
\VUgras{அகல வீதி} \textsuperscript{(19)} \VUgras{சேமிப்பு வங்கிக்குப்}
\textsuperscript{(18)} பின்னால் கடந்து, \VUgras{பெருந் தேவாலயத்தின்}
\textsuperscript{(20)} சதுக்கத்தில் முடிகிறது.

%%K t11-c1-07-1 p
7. --- இந்த நினைவுச் சின்னங்கள் எல்லாம் ஒரு சிறு குன்றின் மேல்
அரங்கம் போல விரிகின்றன ; அங்கே நம் அகன்ற \VUgras{மேல்நிலைப் பள்ளியும்}
\textsuperscript{(21)} இருக்கிறது.

%%K t11-c1-07-2 p
கொஞ்சம் சாய்வான ஒரு வழியில் இறங்கினால் — அது பெரிய தெருவுடன்
இணைகிறது — \VUgras{நீதிமன்றத்தை} \textsuperscript{(22)}
அடைகிறோம் ; அதன் முன்னால் இனிமையான \VUgras{பொதுத் தோட்டம்}
\textsuperscript{(26)} விரிகிறது. இந்த \VUgras{தோட்டச் சதுக்கம்}
மிகப் பெரியது அல்ல ; ஆனால் நகரத்தின் நடுவில் பசுமையின், குளிர்ச்சியின்
சோலையை அது உண்டாக்குகிறது. அதனால் நகரத்தார் அங்கே மிகுதியாக
வருகிறார்கள் ; \VUgras{காப்பி விடுதியின்} \textsuperscript{(25)}
கூடாரத்தின் கீழ் வந்து அமர்வதை அவர்கள் விரும்புகிறார்கள் ; வியாழனும்
ஞாயிறும் \VUgras{இசைக் கூடத்தில்} \textsuperscript{(27)} வாசிக்கும்
படை இசைக் கலைஞர்களைக் கேட்பதற்காக ; அந்தக் கூடம் புல்வெளிகளுக்கும்
செடித் திரள்களுக்கும் இடையே அமைந்திருக்கிறது.

%%K t11-c1-08-1 p
8. --- அந்தப் புல்வெளிகளுள் ஒன்றின் மேல் வெண்கல \VUgras{சிலை}
\textsuperscript{(28)} நிற்கிறது ; அது நம் நகரத்தாருள் ஒருவரின்
நினைவாக எழுப்பப்பட்ட நினைவுச் சின்னம் ; அவர் தான் பிறந்த நகரத்திற்குப்
பெரிய தொண்டுகள் செய்தார்.

%%K t11-tit-3 sub
\VUsaut{4.0mm}
\VUpk{10.2pt}{40}{போக்குவரத்துச் சாதனங்கள்.}
\VUsaut{3.4mm}

%%K t11-c1-09-1 p
9. --- இதுவரை நம் நகரத்திற்கு மின் விளக்கு ஒளி இல்லை ; \VUgras{வாயு
விளக்குகள்} \textsuperscript{(29)} மட்டுமே இருக்கின்றன. ஆனால்
\VUgras{இங்குள்ள செய்தித்தாள்கள்} இந்த ஒளி முறையை மேம்படுத்தப்
பரப்புரை செய்கின்றன ; \VUgras{தண்டவாள வண்டிகளின் இழுவை முறையை}
ஏற்கெனவே மேம்படுத்தியது \VUgras{போல}. \VUgras{குதிரைகள்}
\VUgras{இழுத்த} \VUgras{பழைய} வண்டிகள் நடைமுறையில் \VUgras{மின்
தண்டவாள வண்டிகளால்} \textsuperscript{(31)} மாற்றப்பட்டன ; அவை
மிகவும் \VUgras{மலிவான} விலைக்கு \VUgras{நகரத்தின்} ஒரு நுனியிலிருந்து
\VUgras{மறு நுனிக்கு} பயணிகளை விரைவாகக் கொண்டு செல்கின்றன.

%%K t11-c1-10-1 p
10. --- நீதிமன்றத்திற்கு அருகில் \VUgras{வங்கி} \textsuperscript{(32)}
இருக்கிறது ; அங்கே வணிக உண்டியல்களைத் தள்ளுபடி செய்கிறார்கள்.
\VUgras{வங்கி முகவர்கள்} \textsuperscript{(33)} கடன்களை வீடு தேடி
வசூலிக்கப் போகிறார்கள்.

%%K t11-tit-4 sub
\VUsaut{4.0mm}
\VUpk{10.2pt}{40}{கலையும் கல்வியும்.}
\VUsaut{3.4mm}

%%K t11-c1-11-1 p
11. --- அருகில் நிற்கும் அழகான கட்டிடம் \VUgras{அருங்காட்சியகம்}.
அதில் நகரத்தின் \VUgras{நூலகமும்} \textsuperscript{(34)}, அழகான
\VUgras{இயற்கை அறிவியல் சேகரிப்புகளும்} \textsuperscript{(35)}
(இயற்கை அருங்காட்சியகம்), நேர்த்தியான \VUgras{ஓவிய
அருங்காட்சியகமும்} \textsuperscript{(36)} சேர்ந்திருக்கின்றன ; அந்த
ஓவிய அருங்காட்சியகம் ஒரு சிறந்த நிரந்தர \VUgras{காட்சியகத்தை}
உண்டாக்குகிறது.

%%K t11-c1-11-2 p
வாரத்திற்கு இரண்டு முறை அது பொது மக்களுக்குத் திறக்கப்படுகிறது ;
ஆனால் அயல்நாட்டவர் \VUgras{காவலருக்கு} \textsuperscript{(37)} சிறு
தொகை கொடுத்து எந்த நாளும் பார்க்கலாம். \VUgras{ஓவியர்கள்}
\textsuperscript{(38)} அங்கே வேலை செய்ய வருகிறார்கள். தங்கள்
முக்காலியின் \VUgras{முன்னால்}, கையில் வண்ணப் \VUgras{பலகையுடன்},
புகழ் பெற்ற ஓவியங்களை நகல் எடுப்பதை அவர்களிடம் பார்க்கிறோம்.

%%K t11-c1-12-1 p
12. --- அருங்காட்சியகத்திற்குப் பின்னால், மேட்டின் மேல்,
\VUgras{குவிமாடம்} \textsuperscript{(40)} தெரியத் தொடங்குகிறது ; அது
\VUgras{மருத்துவமனையினுடையது} \textsuperscript{(39)}. \VUgras{ஆசிரியர்
பயிற்சிப் பள்ளியின்} \textsuperscript{(41)} கட்டிடங்களும் தெரிகின்றன ;
அங்கே நம் ஊராட்சிப் பள்ளிகளின் ஆசிரியர்கள் பயிற்சி பெற்றார்கள்.

நாடக அரங்கின் சதுக்கத்தில் \VUgras{அஞ்சல் அலுவலகம்}
\textsuperscript{(43)} இருக்கிறது ; அது ஒரு \VUgras{தனி வீட்டில்}
\textsuperscript{(42)} அமைந்திருக்கிறது.

%%K t11-tit-5 sub
\VUsaut{5.0mm}
\VUpk{11.4pt}{40}{தீ விபத்து.}
\VUsaut{3.0mm}

%%K t11-tit-6 sub
\VUpk{10.2pt}{40}{தீ! என்னும் கூக்குரல்.}
\VUsaut{3.4mm}

%%K t11-c2-01-1 p
1. --- அச்சுறுத்தும் செய்தி நகரம் முழுவதும் பரவுகிறது : நாடக அரங்கில்
இப்போதுதான் தீ விபத்து நேர்ந்தது! நான் \VUgras{சதுக்கத்திற்கு}
\textsuperscript{(96)} வந்த கணத்திலேயே கூட்டம் \VUgras{நடைபாதையிலும்}
\textsuperscript{(46)} \VUgras{வண்டிப் பாதையிலும்} \textsuperscript{(47)}
கூடி விட்டது. \VUgras{அஞ்சல் ஊழியர்களும்} \textsuperscript{(44)}
\VUgras{கடிதம் விநியோகிப்பவர்களும்} \textsuperscript{(45)} அஞ்சல்
அலுவலகத்திலிருந்து வெளியே வருகிறார்கள். \VUgras{தண்டவாள வண்டி
ஓட்டுநர்} \textsuperscript{(48)} தன் வண்டியை நிறுத்த வேண்டியிருந்தது ;
நகரத்தின் \VUgras{நோயாளர் வண்டியை} \textsuperscript{(50)} கடந்து போக
விடுவதற்காக ; அது \VUgras{அறுவை மருத்துவருடனும்} \textsuperscript{(52)}
\VUgras{செவிலியருடனும்} \textsuperscript{(51)} வருகிறது —
\VUgras{காயமுற்றவரைக்} \textsuperscript{(53)} கவனிப்பதற்காக ; அவரை
\VUgras{தூக்குப் படுக்கையில்} \textsuperscript{(54)} \VUgras{செய்தித்தாள்
கடைக்கு} \textsuperscript{(55)} அருகில் கொண்டு வந்தார்கள்.

%%K t11-c2-02-1 p
2. --- பயிற்சியிலிருந்து திரும்பி வந்து கொண்டிருந்த காலாட்படைப் பிரிவு
ஒன்று நின்று விட்டது ; \VUgras{குதிரைப் படைக் காவலர்களுடன்}
\textsuperscript{(61)} சேர்ந்து ஒழுங்கைக் காப்பதற்காக.
\VUgras{குதிரை வீரர்கள்} — அவர்களுடைய குதிரைகள் பின்னங்கால்களில்
எழுகின்றன — \VUgras{வீரர்களை} \textsuperscript{(57)} விடவும்
\VUgras{படைக் காவலர்களை} \textsuperscript{(63)} விடவும் நன்றாக
\VUgras{வேடிக்கை பார்ப்பவர்களைப்} \textsuperscript{(56)} பின்னுக்குத்
தள்ளுகிறார்கள். மேலும் படைக் காவலர்கள் மிகவும் பணியில்
இருக்கிறார்கள். அவர்களுள் ஒருவர் \VUgras{காவல் துறையினருக்கு}
\textsuperscript{(64)} இன்னொரு \VUgras{விபத்துக்கு உள்ளானவரை}
\textsuperscript{(65)} எங்கே கொண்டு போக வேண்டும் என்று
காட்டுகிறார் ; அவர் ஏணியிலிருந்து விழுந்த துணிச்சலான தீயணைப்பு
வீரர்.

%%K t11-c2-03-1 p
3. --- \VUgras{அதிகாரி} \textsuperscript{(66)} வரவிருக்கும்
\VUgras{நீராவி நீர்ப் பீச்சியை} \textsuperscript{(67)} எங்கே வைக்க
வேண்டும் என்பதைத் துல்லியமாகச் சொல்கிறார். அந்த வலிமையான
இயந்திரத்திற்குக் காத்திருக்கும் போது, அவசரமாக ஒரு \VUgras{கை
நீர்ப் பீச்சியைப்} \textsuperscript{(68)} பொருத்த வைத்தார் ; அதன்
நீளமான \VUgras{குழாய்} \textsuperscript{(69)} தீயின் மேல் நீரை
வெள்ளமாக வீசுகிறது.

ஆறு தீயணைப்பு வீரர்கள் அதைச் செலுத்துகிறார்கள் ; அதே வேளையில்
அவர்களுடைய தோழன் \VUgras{பீச்சு முனையைப்} \textsuperscript{(70)}
பிடித்து \VUgras{நீர்ப் பீச்சை} \textsuperscript{(71)} தீ விபத்தின்
மையத்தை நோக்கிச் செலுத்துகிறான்.

%%K t11-c2-04-1 p
4. --- நாடக அரங்கின் மறு பக்கத்தில் பெரிய \VUgras{ஏணியைச்}
\textsuperscript{(72)} சாய்த்து வைத்தார்கள். கருமையான \VUgras{புகைச்}
\textsuperscript{(75)} சுழல்களுக்கு இடையே பீறிடும் தீச் சுவாலைகளை
அணுக அது தீயணைப்பு வீரர்களுக்கு வழி செய்கிறது.

%%K t11-tit-7 sub
\VUsaut{5.0mm}
\VUpk{10.2pt}{40}{துணிச்சலான செயல்கள்.}
\VUsaut{3.4mm}

%%K t11-c2-05-1 p
5. --- \VUgras{தீ விபத்து} \textsuperscript{(74)} \VUgras{நாடக
அரங்கின்} \textsuperscript{(73)} முன் கூடத்தில் பிறந்தது ; அப்போது
\VUgras{இசை நாடகம்} ஒன்றை ஒத்திகை பார்த்துக் கொண்டிருந்தார்கள் ;
அதன் முதல் \VUgras{காட்சி} மறு நாள் நடந்திருக்கும். நல்ல வேளையாகக்
காட்சிக் கூடத்தில் சில \VUgras{பார்வையாளர்கள்} \textsuperscript{(76)}
மட்டுமே இருந்தார்கள். அந்தப் பரிதாபமானவர்கள் \VUgras{பால்கனியில்}
\textsuperscript{(77)} அடைக்கலம் புகுந்தார்கள் ; அங்கே துணிச்சலான
\VUgras{தீயணைப்பு வீரர்கள்} \textsuperscript{(86)} ஏணியாலும் பெருங்
கயிறுகளாலும் அவர்களைக் காப்பாற்ற வருகிறார்கள்.

%%K t11-c2-06-1 p
6. --- \VUgras{தூண் வரிசை மண்டபத்தின்} \textsuperscript{(78)} கீழ் —
அங்கே \VUgras{விளம்பரத் தாள்} \textsuperscript{(79)} காட்சியைப் பற்றித்
தெரிவிக்கிறது — \VUgras{இசைக் கலைஞர்கள்} \textsuperscript{(81)}
ஓடுவதைப் பார்க்கிறோம் ; அவர்கள் \VUgras{இசைக் குழுவினர்}
\textsuperscript{(80)} ; தங்கள் இசைக் கருவிகளைச் சுமந்து செல்கிறார்கள் :
\VUgras{வயலின்} \textsuperscript{(81)}, \VUgras{புல்லாங்குழல்}
\textsuperscript{(82)}, \VUgras{செலோ} \textsuperscript{(83)},
\VUgras{ட்ரோம்போன்} \textsuperscript{(84)}, \VUgras{முரசு}
\textsuperscript{(85)} முதலியன. \VUgras{தளவாடங்களில்}
\textsuperscript{(87)} ஒரு பகுதியைக் காப்பாற்ற முயல்கிறார்கள்.

%%K t11-c2-07-1 p
7. --- \VUgras{நடிகர்களும்} \textsuperscript{(89)}
\VUgras{நடிகைகளும்} \textsuperscript{(88)}, \VUgras{பாடகர்களும்}
\VUgras{பாடகிகளும்} தங்கள் நாடக \VUgras{உடையுடன்}
\textsuperscript{(90)} ஓட வேண்டியிருந்தது. உயர்குரல் பாடகர் ஒரு
\VUgras{செய்தியாளருக்கு} \textsuperscript{(92)} இது எப்படி நேர்ந்தது
என்று விளக்குகிறார். \VUgras{நிருபர்} முதல் கணத்திலிருந்தே
அதிகாரிகளுடன் ஓடி வந்தார் : \VUgras{ஆட்சியர்} \textsuperscript{(91)},
\VUgras{மேயர்} \textsuperscript{(93)}, \VUgras{காவல் ஆணையர்}
\textsuperscript{(94)}, \VUgras{நீதித் துறை அலுவலர்}
\textsuperscript{(95)} ; அவர்கள் பொறுப்புகளைத் தீர்ப்பதற்காக விசாரணை
செய்வார்கள்.

\VUsaut{6.0mm}
\VUfilet{22mm}
